% Section 5: Zero-shot context selection with controllable risk
% Required packages: \usepackage{amsmath,booktabs,graphicx}
% Upload this file together with the three PDF figures under figures/.

\section{Zero-shot context selection with controllable risk}
\label{sec:zero-shot-predictor}

Sections~\ref{sec:real-context}--\ref{sec:context-mechanisms} established that
the native context is often unnecessarily long, but also that the best shorter
window is not stable enough to treat as an ordinary classification label.  The
deployment question is therefore not simply \emph{which window will win?}  It
is \emph{how much risk of degradation are we willing to accept in exchange for
less context compute?}

We answer that question with a zero-shot risk predictor.  For every candidate
window, it estimates loss relative to native context, uncertainty in that
estimate, and the probability of a material error increase.  Synthetic-only
calibration converts these estimates into five ordered operating profiles,
from Conservative to Max efficiency.  A profile selects the shortest window
that passes its risk threshold and otherwise abstains to native context.  The
result is an explicit compute--harm dial rather than a single brittle guess.

\subsection{Predict risk relative to the native action}
\label{sec:risk-target}

Each frozen TSFM labels a synthetic pool independently over six forecast
horizons, $h\in\{16,32,64,128,512,1024\}$, and an 18-point context grid from
32 to 8,192 steps, truncated to the model's supported context.  For synthetic
series $i$, candidate window $L_k$, horizon $h$, and forecast MAE
$E_{i,k,h}$, the regression target is the log-risk relative to the model's
native window,

\begin{equation}
    r_{i,k,h}
      = \log\!\left(\frac{E_{i,k,h}}
                          {E_{i,\mathrm{native},h}}\right).
    \label{eq:risk-target}
\end{equation}

The regressor is fitted to this target after clipping it to $[-3,3]$.  A
separate classifier predicts a material-harm event,

\begin{equation}
    b_{i,k,h}
      = \mathbb{1}\!\left[
          E_{i,k,h} > 1.05\,E_{i,\mathrm{native},h}
        \right].
    \label{eq:harm-event}
\end{equation}

This relative formulation makes native context the explicit safe action and
avoids asking the predictor to reproduce a noisy exact argmin.  Positive risk
means that the shorter window is expected to hurt; negative risk means that it
is expected to improve on native context.

The predictor is deliberately lightweight.  A 192-tree ExtraTrees regressor
and classifier consume 94 scale-invariant features describing recent and
remote summaries, autocorrelation, spectrum, regime shifts, valid-history
length, forecast horizon, and candidate-window geometry.  For each TSFM, the
models use 12,000 synthetic training series, 2,000 disjoint synthetic
calibration series, and 600,000 sampled series--window--horizon pairs.  No
GiftEval target, error, or model-selection signal enters fitting or
calibration.

\begin{figure*}[t]
    \centering
    \includegraphics[width=\textwidth]{figures/fig1_risk_selector_pipeline.pdf}
    \caption{Risk-calibrated context selection.  Offline, a frozen TSFM labels
    synthetic context actions relative to its native window; a lightweight
    predictor learns expected log-risk, uncertainty, and the probability of
    more than 5\% harm.  Synthetic calibration turns one frozen risk score into
    five nested profiles.  At deployment, the chosen profile returns the
    shortest acceptable window or abstains to native context.  Real forecast
    outcomes are revealed only for final evaluation.}
    \label{fig:risk-pipeline}
\end{figure*}

\subsection{Calibrate a selectable harm profile}
\label{sec:risk-calibration}

Let $\widehat\mu_{i,k,h}$ be predicted mean log-risk,
$\widehat\sigma_{i,k,h}$ the standard deviation across regression trees, and
$\widehat p^{(5)}_{i,k,h}$ the predicted probability of the event in
Equation~\ref{eq:harm-event}.  Candidate actions are ordered by the composite
risk score

\begin{equation}
    q_{i,k,h}
      = \widehat\mu_{i,k,h}
        + u\,\widehat\sigma_{i,k,h}
        + v\,\widehat p^{(5)}_{i,k,h},
    \label{eq:risk-score}
\end{equation}

where $(u,v)$ is selected once for each TSFM using the balanced synthetic
profile.  Holding that score fixed, profile-specific thresholds are calibrated
on the same disjoint synthetic split.  The deployed decision is

\begin{equation}
    \widehat L_{i,h}^{(a)}
      = \min\{L_k:q_{i,k,h}\leq\tau_a\},
    \quad\text{or }L_{\mathrm{native}}\text{ if the set is empty},
    \label{eq:risk-selection}
\end{equation}

for profile $a$.  For Chronos2-Base, calibration selects $(u,v)=(0,0.5)$, so
the deployed score is the predicted log-error ratio plus one half of the
predicted probability of more than 5\% harm.  This makes the safety preference
explicit without changing the shortest-acceptable-window decision rule.

Because a looser profile raises the accepted-risk threshold,
the profiles are nested: increasing the harm budget can only retain or shorten
the chosen context.  Calibration constrains both the rate of instances harmed
by more than 5\% and the mean error ratio.  The five budgets and their real
zero-shot operating points are reported together in
Table~\ref{tab:risk-profiles}.

\begin{table*}[t]
    \centering
    \footnotesize
    \renewcommand{\arraystretch}{1.12}
    \caption{The calibrated compute--harm dial.  Budget columns are constraints
    on the disjoint synthetic calibration split.  Real columns aggregate 97
    GiftEval cells for each of 11 TSFMs: coverage and FLOPs saved are macro
    averages across models, harm is pooled across model--instance evaluations,
    and $\Delta$MASE is the change implied by the geometric mean of the 11
    cell-geometric MASE ratios.  Positive $\Delta$MASE denotes improvement.}
    \label{tab:risk-profiles}
    \begin{tabular*}{\textwidth}{@{\extracolsep{\fill}}lrrrrrr@{}}
        \toprule
        Profile & Harm budget & Mean-ratio cap & Coverage & FLOPs saved & Real harm $>5\%$ & $\Delta$MASE \\
        \midrule
        Conservative   & 0.5\% & 1.000 &  6.43\% &  4.77\% & 0.70\% & $-0.015\%$ \\
        Balanced       & 1.0\% & 1.001 & 17.77\% &  9.04\% & 1.54\% & $-0.014\%$ \\
        Aggressive     & 3.0\% & 1.005 & 29.77\% & 19.17\% & 3.01\% & $-0.045\%$ \\
        \addlinespace[2pt]
        Efficiency      & 15.0\% & 1.040 & 49.21\% & 45.43\% & 8.52\% & $-0.401\%$ \\
        Max efficiency  & 20.0\% & 1.070 & 56.01\% & 52.57\% & 11.15\% & $-0.795\%$ \\
        \bottomrule
    \end{tabular*}
\end{table*}

The central result is the ordering of these operating points.  Moving from the
Conservative to the Max-efficiency profile increases average theoretical
TSFM FLOPs savings from 4.77\% to 52.57\%, while the pooled rate of instances
whose MASE worsens by more than 5\% rises from 0.70\% to 11.15\%.  The
corresponding model-geometric MASE change is 0.795\% worse than native at the
most aggressive endpoint.  The risk predictor therefore exposes a useful continuum:
deployments can preserve native behavior almost everywhere, or deliberately
accept more tail risk to unlock substantially more context reduction.

\begin{figure}[t]
    \centering
    \includegraphics[width=\columnwidth]{figures/fig2_compute_harm_dial.pdf}
    \caption{Aggregate real-data behavior of the five synthetic-calibrated
    profiles.  Left: mean theoretical TSFM FLOPs saved versus pooled incidence
    of more than 5\% harm.  Right: aggregate MASE stays close to native while
    profiles move along the compute--harm curve.  The labels show synthetic
    harm budgets; they are calibration targets, not real-domain guarantees.}
    \label{fig:compute-harm-dial}
\end{figure}

\subsection{The same dial behaves differently across forecasters}
\label{sec:risk-transfer}

The average curve hides substantial model dependence
(Table~\ref{tab:risk-models} and Figure~\ref{fig:risk-models}).  At the
Max-efficiency profile, Chronos2-Base and Chronos2-Synth save 77.44\% and
74.18\% of theoretical FLOPs with 12.45\% and 15.48\% observed harm,
respectively.  Chronos2-Small saves 64.73\% with a 0.563\% aggregate MASE cost.
Moirai2-Small reaches 70.13\% savings at 11.16\% harm.  ChronosBolt-Base is
much more cautious, saving 20.71\% with a 0.070\% aggregate MASE cost.

The Chronos2-Base result is also meaningful across model sizes.  Under Max
efficiency, its cell-geometric MASE is 0.9995, compared with 1.0176 for
full-context Chronos2-Small, while its selected theoretical compute is 77.4\%
below the full-context Small run.  On this benchmark, the risk-controlled Base
model therefore beats full Small in both aggregate accuracy and theoretical
compute, making Max efficiency the strongest efficiency endpoint rather than the
safest operating default.

Other models are naturally conservative.  Toto and ChronosBolt abstain almost
entirely under the conservative profile.  Sundial shortens many instances but
realizes little theoretical FLOPs reduction because many accepted changes are
small relative to its native compute profile.  These cases are not failures of
the abstention mechanism: they show that the same harm target can imply very
different amounts of useful compression.

\begin{table*}[t]
    \centering
    \footnotesize
    \renewcommand{\arraystretch}{1.10}
    \caption{Per-forecaster operating points on GiftEval.  The balanced and
    Max-efficiency columns report theoretical TSFM FLOPs saved and the
    instance rate with more than 5\% MASE harm.  The final column reports the
    cell-geometric MASE change under the Max-efficiency profile; positive is
    better.}
    \label{tab:risk-models}
    \begin{tabular*}{\textwidth}{@{\extracolsep{\fill}}lrrrrr@{}}
        \toprule
        Model & Balanced FLOPs & Balanced harm & Max-eff. FLOPs & Max-eff. harm & $\Delta$MASE \\
        \midrule
        Chronos2-Base      & 18.05\% & 0.96\% & 77.44\% & 12.45\% & $-1.368\%$ \\
        Chronos2-Small     & 12.44\% & 2.13\% & 64.73\% &  9.20\% & $-0.563\%$ \\
        Chronos2-Synth     & 15.05\% & 0.96\% & 74.18\% & 15.48\% & $-1.301\%$ \\
        ChronosBolt-Base   &  1.95\% & 0.39\% & 20.71\% &  7.89\% & $-0.070\%$ \\
        FlowState-R1       & 18.90\% & 5.58\% & 48.30\% & 10.30\% & $-0.844\%$ \\
        Moirai2-Small      & 12.15\% & 1.00\% & 70.13\% & 11.16\% & $-0.352\%$ \\
        PatchTST-FM-R1     &  8.45\% & 1.05\% & 49.89\% &  7.58\% & $-0.747\%$ \\
        Sundial-Base-128M  &  0.72\% & 3.85\% &  4.93\% & 16.15\% & $-0.455\%$ \\
        TimesFM2.5-200M    &  6.87\% & 0.39\% & 70.62\% & 11.27\% & $-0.900\%$ \\
        TiRex2             &  4.82\% & 0.45\% & 53.34\% & 11.44\% & $-1.243\%$ \\
        Toto-2.0-313m      &  0.01\% & 0.01\% & 44.02\% &  9.77\% & $-0.914\%$ \\
        \bottomrule
    \end{tabular*}
\end{table*}

\begin{figure*}[t]
    \centering
    \includegraphics[width=\textwidth]{figures/fig3_model_profile_heatmaps.pdf}
    \caption{How the risk dial transfers across forecasters.  Each row is one
    independently trained predictor and each column is a harm profile.  Cells
    show theoretical FLOPs saved (left) and observed incidence of more than
    5\% MASE harm (right).  The nested profiles produce an ordered trade-off,
    but the attainable savings and calibration error remain model specific.}
    \label{fig:risk-models}
\end{figure*}

\subsection{Where the efficiency profiles help---and where they fail}
\label{sec:risk-datasets}

The model-wide result is not uniform across recognizable data domains.
Table~\ref{tab:risk-datasets} groups the Chronos2-Base cells by dataset family.
Here $\Delta$MASE is relative to full-context Chronos2-Base (negative is
better), context saved pools selected and native context tokens across all
family instances, and shortened is the fraction that does not abstain.

\begin{table*}[t]
    \centering
    \footnotesize
    \renewcommand{\arraystretch}{1.12}
    \caption{Dataset-level behavior of the Chronos2-Base Max-efficiency profile.
    Large aggregate context savings can coexist with high instance-level tail
    risk.  M4 and COVID Deaths show the opposite regime: selective or complete
    abstention.}
    \label{tab:risk-datasets}
    \begin{tabular*}{\textwidth}{@{\extracolsep{\fill}}lrrrr@{}}
        \toprule
        Dataset family & $\Delta$MASE & Context saved & Shortened & Harm $>5\%$ \\
        \midrule
        M4              & $-1.12\%$ & 11.5\% &  5.1\% &  1.4\% \\
        Electricity     & $+0.22\%$ & 80.7\% & 97.0\% & 18.1\% \\
        Solar           & $+0.64\%$ & 72.1\% & 95.4\% & 23.2\% \\
        Shenzhen Taxi   & $+0.79\%$ & 69.0\% & 99.1\% & 17.8\% \\
        ETTm1           & $+1.01\%$ & 75.3\% & 97.0\% & 27.5\% \\
        Jena Weather    & $+1.89\%$ & 76.7\% & 97.1\% & 35.6\% \\
        ETTm2           & $+2.06\%$ & 77.3\% & 95.9\% & 25.8\% \\
        Seattle Traffic & $+2.39\%$ & 76.1\% & 97.5\% & 30.8\% \\
        \addlinespace[2pt]
        COVID Deaths    & $+0.00\%$ &  0.0\% &  0.0\% &  0.0\% \\
        \bottomrule
    \end{tabular*}
\end{table*}

M4 is the cleanest reliability result.  Short histories make the policy abstain
on 94.9\% of instances; the remaining selective shortening reduces aggregate
MASE by 1.12\%.  COVID Deaths triggers complete abstention, the intended
response when no shorter action appears acceptable.  Electricity is the
strongest efficiency result: 80.7\% of context is removed for only a 0.22\%
aggregate MASE cost.  Solar is also attractive in aggregate, with 72.1\%
context reduction for a 0.64\% MASE cost.

The same aggregates expose a reliability gap.  ETT and Jena Weather are
horizon sensitive: some long- or medium-horizon cells improve while their
short-horizon counterparts degrade.  For example, ETTm1 hourly-medium improves
by about 2\%, whereas ETTm1 hourly-short worsens by about 5.3\%.  Jena Weather
10-minute long improves by 2.6\%, but its short-horizon cell worsens by 9.8\%.
Seattle Traffic is the weakest familiar family, combining 76.1\% context
reduction with a 2.39\% aggregate MASE increase and 30.8\% instance-level harm.

The Efficiency profile is a more credible operational choice when individual
series reliability matters.  It retains most of the compute benefit on
Electricity (75.9\% context saved) while improving MASE by 0.27\%.  Relative to
Max efficiency, it also lowers MASE costs on ETTm1 (0.22\% versus 1.01\%), ETTm2
(1.04\% versus 2.06\%), Jena Weather (1.67\% versus 1.89\%), and Seattle
Traffic (1.73\% versus 2.39\%).  Solar is the exception: Max efficiency has the
lower aggregate MASE cost (0.64\% versus 1.22\%) despite its looser threshold.
M4 remains
selective and improves by 1.47\%.  Thus Max efficiency is the efficiency-first
endpoint, whereas Efficiency is the stronger default for reliability-sensitive
deployment.

These discrete profiles are operating-point samples, not the final frontier.
We will next evaluate the dense threshold sweep and retain only nondominated
accuracy--compute points to construct a Pareto front.  That analysis will show
whether the named profiles lie near the attainable frontier or whether better
thresholds dominate them.

\subsection{Calibration is a control target, not a guarantee}
\label{sec:risk-limits}

The profiles are calibrated without real labels, so their budgets should be
read as deployment controls rather than distribution-free guarantees.  This
distinction is visible in the balanced profile: its synthetic harm budget is
1\%, whereas pooled real harm is 1.54\%.  FlowState-R1 is the clearest
calibration failure.  Even its conservative profile produces 5.58\% real harm
and a 0.330\% aggregate MASE degradation, despite satisfying the 0.5\%
synthetic calibration target.  For this model, native context remains the
appropriate default until calibration data better matches the deployment
distribution.

Two compute qualifications also remain.  Reported savings cover the TSFM
forecast stage and do not yet include feature extraction or tree inference.
They are theoretical proxies: TiRex2 internally pads inputs to a fixed length,
and the quadratic proxy is only approximate for the linear-state-space
FlowState architecture.  These two models therefore require measured latency
and memory profiling before their apparent context savings are interpreted as
deployment speedups.

\subsection{Implications}
\label{sec:risk-implications}

The useful-context problem is better posed as selective risk management than
as exact window prediction.  Relative-risk targets align learning with the
native action, abstention prevents forced shortening, and calibrated profiles
make the deployment preference explicit.  Across 11 forecasters, this design
turns one zero-shot predictor into a family of operating points spanning
roughly 5--53\% average theoretical FLOPs savings.  Its main limitation is
equally informative: synthetic calibration does not transfer uniformly across
architectures.  The practical next step is therefore not a more aggressive
argmin predictor, but stronger calibration diagnostics and model-specific
safeguards around a user-chosen harm budget.
